\documentclass[12pt,a4paper]{article}
\usepackage[utf8]{inputenc}
\usepackage[T1]{fontenc}
\usepackage[spanish]{babel}
\usepackage{geometry}
\geometry{margin=2.5cm}
\usepackage{hyperref}
\usepackage{booktabs}
\usepackage{array}
\usepackage{enumitem}
\usepackage{amsmath}
\usepackage{fancyhdr}

\pagestyle{fancy}
\fancyhf{}
\rhead{\small Actividad 1 -- Semana 1}
\lhead{\small Sistemas de Información -- UNIDEP}
\cfoot{\thepage}

\title{\textbf{Documentación Técnica, Marco Teórico y Evaluación}\\
\large Construcción de una API REST con Flask}
\author{Profesor: Sergio Mirazo \\ Materia: Sistemas de Información\\
Ingeniería en Sistemas Computacionales -- UNIDEP}
\date{}

\begin{document}
\maketitle
\thispagestyle{empty}

\section*{Documentación de Endpoints y Estructura del Proyecto}

\subsection*{Tabla de Endpoints}
\begin{table}[h]
\centering
\renewcommand{\arraystretch}{1.3}
\begin{tabular}{|c|l|p{6cm}|c|}
\hline
\textbf{Método} & \textbf{Ruta} & \textbf{Descripción} & \textbf{Códigos HTTP} \\ \hline
GET    & \verb|/api/recursos|          & Retorna todos los registros almacenados.        & 200 \\ \hline
GET    & \verb|/api/recursos/<id>|     & Retorna un recurso específico por su identificador. & 200, 404 \\ \hline
POST   & \verb|/api/recursos|          & Crea un nuevo registro. Requiere cuerpo JSON.   & 201, 400 \\ \hline
PUT    & \verb|/api/recursos/<id>|     & Actualiza parcialmente o completamente un registro existente. & 200, 400, 404 \\ \hline
DELETE & \verb|/api/recursos/<id>|     & Elimina un registro por su identificador.       & 200, 404 \\ \hline
\end{tabular}
\end{table}

\subsection*{Ejemplo de Petición y Respuesta (POST)}
\textbf{Cuerpo de la petición:}
\begin{verbatim}
{
  "nombre": "Pedro Sánchez",
  "email": "pedro@email.com",
  "rol": "DevOps"
}
\end{verbatim}

\textbf{Respuesta esperada (201 Created):}
\begin{verbatim}
{
  "mensaje": "Recurso creado",
  "id": 4,
  "dato": {
    "nombre": "Pedro Sánchez",
    "email": "pedro@email.com",
    "rol": "DevOps"
  }
}
\end{verbatim}

\subsection*{Estructura del Proyecto}
\begin{verbatim}
api_unidep_semana1/
|-- main.py          # Lógica de la API, rutas y controladores
|-- bd.py            # Almacén simulado en memoria (diccionario)
`-- README.md        # Documentación técnica y guía de ejecución
\end{verbatim}

\textbf{Notas técnicas:} La base de datos se mantiene en volátil (memoria RAM del proceso). Al reiniciar el servidor, los datos vuelven a su estado inicial definido en \texttt{bd.py}. Este enfoque es válido para etapas de prototipado y aprendizaje de flujo HTTP.

\section*{Marco Teórico de Reforzamiento}

\subsection*{1. Fundamentos de una API REST}
Una \textbf{API (Application Programming Interface)} es un conjunto de reglas y protocolos que permite la comunicación entre sistemas de software. El estilo \textbf{REST (Representational State Transfer)} se basa en principios arquitectónicos definidos por Roy Fielding:
\begin{itemize}
    \item \textbf{Cliente-Servidor:} Separación de responsabilidades para mejorar portabilidad y escalabilidad.
    \item \textbf{Stateless (Sin estado):} Cada petición contiene toda la información necesaria para ser procesada. El servidor no guarda contexto de sesión entre peticiones.
    \item \textbf{Interfaz uniforme:} Uso estandarizado de métodos HTTP, URI identificadores y formatos de representación (JSON, XML).
    \item \textbf{Cacheable:} Las respuestas pueden marcarse como cacheables para reducir latencia y carga.
\end{itemize}

\subsection*{2. Protocolo HTTP y Códigos de Estado}
HTTP (Hypertext Transfer Protocol) es el protocolo subyacente de las APIs REST. Su estructura se compone de:
\begin{itemize}
    \item \textbf{Línea de petición:} Método, URI y versión HTTP.
    \item \textbf{Encabezados (Headers):} Metadatos como \texttt{Content-Type}, \texttt{Authorization}, \texttt{Accept}.
    \item \textbf{Cuerpo (Body):} Payload de la petición o respuesta (comúnmente JSON).
\end{itemize}

\textbf{Códigos de estado relevantes:}
\begin{itemize}
    \item \textbf{2xx (Éxito):} 200 OK, 201 Created, 204 No Content.
    \item \textbf{4xx (Error del cliente):} 400 Bad Request, 401 Unauthorized, 403 Forbidden, 404 Not Found.
    \item \textbf{5xx (Error del servidor):} 500 Internal Server Error, 502 Bad Gateway.
\end{itemize}

\subsection*{3. Arquitectura: Data Warehouse vs Data Lakehouse}
\begin{itemize}
    \item \textbf{Data Warehouse (DW):} Repositorio centralizado de datos estructurados y procesados. Utiliza el paradigma \textit{Schema-on-Write}, optimizado para consultas analíticas (OLAP) y reportes de BI. Su gobernanza es estricta, pero su escalabilidad horizontal es costosa.
    \item \textbf{Data Lake:} Almacén de datos crudos en cualquier formato (estructurados, semiestructurados, no estructurados). Usa \textit{Schema-on-Read}, ofreciendo flexibilidad y bajo costo, pero con desafíos en calidad, gobernanza y consistencia ACID.
    \item \textbf{Data Lakehouse:} Arquitectura híbrida que combina la flexibilidad y economía del Data Lake con la gestión transaccional, gobernanza y rendimiento del DW. Utiliza formatos abiertos como Delta Lake, Apache Iceberg o Hudi para soportar transacciones ACID, time-travel y consumo directo desde motores de BI y Machine Learning sin necesidad de duplicar datos.
\end{itemize}

\subsection*{4. Pipelines para APIs}
Las APIs modernas son componentes críticos en los \textbf{pipelines de datos}:
\begin{itemize}
    \item \textbf{Ingestión:} APIs públicas o privadas extraen datos de fuentes externas (SaaS, sensores, logs).
    \item \textbf{Transformación y Orquestación:} Herramientas como Apache Airflow, Prefect o dbt consumen endpoints, validan esquemas y aplican reglas de negocio.
    \item \textbf{API Gateways:} Puntos de entrada que gestionan autenticación, rate limiting, caché, enrutamiento y monitoreo antes de llegar a los microservicios.
    \item \textbf{Patrones de entrega:} Síncrono (REST/GraphQL) vs Asíncrono (Webhooks, colas de mensajes como RabbitMQ o Kafka). En arquitecturas de datos modernas, se prioriza el desacoplamiento mediante eventos para garantizar resiliencia y escalabilidad.
\end{itemize}

\subsection*{5. Bases de Datos: SQL vs NoSQL}
\begin{itemize}
    \item \textbf{SQL (Relacionales):} Basadas en tablas con esquemas rígidos. Garantizan propiedades \textbf{ACID} (Atomicidad, Consistencia, Aislamiento, Durabilidad). Ideales para transacciones financieras, relaciones complejas y consistencia estricta. Ejemplos: PostgreSQL, MySQL, Oracle.
    \item \textbf{NoSQL (No Relacionales):} Modelos flexibles: documentos (MongoDB), clave-valor (Redis), columnares (Cassandra) o grafos (Neo4j). Siguen principios \textbf{BASE} (Basically Available, Soft state, Eventual consistency). Optimizadas para escalabilidad horizontal, grandes volúmenes de datos semiestructurados y baja latencia.
    \item \textbf{Criterios de selección:} Volumen y variedad de datos, necesidad de transacciones estrictas, patrones de consulta, requisitos de escalabilidad y tolerancia a fallos. En la práctica, sistemas modernos adoptan arquitecturas \textit{polyglot persistence}, combinando ambos según el caso de uso.
\end{itemize}

\section*{Cuestionario de Reforzamiento}

\begin{enumerate}[label=\arabic*.]
    \item ¿Qué significa que una API REST sea ``stateless'' y cuál es su implicación práctica en el diseño del servidor?
    \item Mencione tres diferencias clave entre los métodos HTTP PUT y PATCH.
    \item ¿Cuál es la función principal de un API Gateway en una arquitectura distribuida?
    \item Explique el concepto de \textit{Schema-on-Write} y en qué arquitectura de datos se aplica tradicionalmente.
    \item ¿Qué ventaja principal ofrece un Data Lakehouse respecto a un Data Warehouse tradicional en entornos de Machine Learning?
    \item Defina el acrónimo ACID y explique por qué es crítico en sistemas transaccionales.
    \item ¿En qué consiste el principio de \textit{Schema-on-Read} y qué tipo de base de datos lo utiliza comúnmente?
    \item Mencione dos escenarios donde una base de datos NoSQL es preferible sobre una relacional.
    \item ¿Qué código HTTP debe retornar una API cuando se intenta actualizar un recurso con un cuerpo JSON malformado?
    \item ¿Cuál es la diferencia fundamental entre un pipeline ETL y un pipeline ELT?
    \item Explique brevemente cómo funciona el mecanismo de rate limiting en un endpoint público.
    \item ¿Qué formato de datos es estándar en las APIs REST modernas y por qué ha reemplazado a XML en muchos casos?
    \item Describa el rol de un webhook en la integración de APIs asíncronas.
    \item ¿Qué propiedad de las bases de datos relacionales garantiza que una transacción se ejecute por completo o no se ejecute en absoluto?
    \item ¿Por qué un Data Lake tradicional puede presentar problemas de gobernanza y calidad de datos a largo plazo?
    \item Mencione dos estrategias para manejar la versión de una API REST (API versioning).
    \item ¿Qué diferencia existe entre escalabilidad vertical y horizontal? ¿Cuál es más común en entornos NoSQL?
    \item Explique por qué el método POST se considera no idempotente, mientras que PUT sí lo es.
    \item ¿Qué rol juegan los índices en el rendimiento de una base de datos y qué costo operativo introducen?
    \item Si una API retorna un error 503 Service Unavailable, ¿qué indica esto sobre el estado del servidor y qué acción debería tomar un cliente resiliente?
\end{enumerate}





\end{document}